\section{Обзор предметной области}
\label{sec:Chapter1} \index{Chapter1}

\subsection{Анализ безопасности OCI-образов}

Современная разработка и эксплуатация программного обеспечения в значительной степени основана на использовании \textit{контейнеров}~--- легковесных, изолированных сред исполнения, обеспечивающих воспроизводимость окружения. На уровне инструментов это направление является развитием идей, изначально сформировавшихся в Docker и впоследствии стандартизированных в спецификациях \textit{Open Container Initiative (OCI)}. OCI определяет три основных раздела стандарта: спецификацию формата образа~\cite{ociimagespec}, описывающую устройство контейнерных образов, спецификацию распространения~\cite{ocidistspec}, описывающую API container registry, а также спецификацию исполнения~\cite{ociruntimespec}, описывающую формат и поведение контейнеров на стороне runtime. Согласованность этих спецификаций обеспечивает взаимозаменяемость инструментов и единый формат артефактов в рамках всего жизненного цикла.

С точки зрения безопасности образ представляет собой обширный артефакт, объединяющий в себе компоненты различной природы. В типичный образ могут входить:

\begin{itemize}
\item базовый Linux-дистрибутив с установленным набором системных пакетов (как правило, через пакетные менеджеры apt, apk, dnf и аналогичные),
\item разделяемые библиотеки и среды исполнения языков программирования (например, glibc, OpenSSL, JVM, интерпретаторы Python, Node.js, Ruby),
\item языковые зависимости приложения (Python-модули, npm-пакеты, Maven-артефакты, Go-модули и т.п.),
\item собственно код приложения, который может находиться как в виде исходных файлов, так и в виде скомпилированных бинарных файлов.
\end{itemize}

Каждая из этих категорий может содержать \textit{известные уязвимости (CVE)}, информация о которых регулярно публикуется в авторитетных базах, в первую очередь в \textit{National Vulnerability Database (NVD)}~\cite{nvd}, поддерживаемой NIST, а также в дистрибутив-специфичных источниках (Debian Security Tracker, Red Hat Security Data API, Alpine SecDB и др.) и в базах языковых экосистем (например, GitHub Advisory Database). Каждой записи в этих базах соответствует уникальный идентификатор формата \texttt{CVE-YYYY-NNNN}, что позволяет однозначно ссылаться на конкретную уязвимость~\cite{cve}.

При наличии уязвимости в каком-либо из компонентов образа уязвимым становится и контейнер, развёрнутый из этого образа. На практике вероятность того, что в произвольно взятом образе обнаружится хотя бы одна известная уязвимость, крайне высока, особенно с учётом темпов пополнения соответствующих баз: на момент написания работы в NVD насчитывается более 150000 записей, и их количество увеличивается на десятки тысяч ежегодно. Регулярная и автоматизированная проверка образов на наличие уязвимостей становится поэтому обязательной частью цикла поставки программного обеспечения и одним из ключевых требований к безопасной эксплуатации контейнерной инфраструктуры.

\subsection{Структура OCI-образа}
\label{subsec:oci-structure}

С технической точки зрения OCI-образ представляет собой структурированный набор объектов, связанных между собой ссылками на основе криптографических хешей содержимого~\cite{ociimagespec}. В состав образа входят:

\begin{itemize}
\item \textit{манифест (manifest)}~--- JSON-документ, описывающий конкретный вариант образа для определенного сочетания операционной системы и архитектуры процессора. В нём перечислены ссылки на конфигурацию образа и на слои файловой системы;
\item \textit{конфигурация (config)}~--- JSON-документ, содержащий параметры запуска контейнера (точку входа, переменные окружения, рабочую директорию и пр.), а также историю команд сборки образа;
\item \textit{слои (layers)}~--- упорядоченный набор бинарных объектов, каждый из которых содержит разностное содержимое файловой системы относительно предыдущего слоя в формате tar-архива, как правило, сжатого алгоритмом gzip;
\item \textit{индекс (image index)}~--- опциональный JSON-документ, объединяющий несколько манифестов в один логический образ. Обычно используется для multi-platform образов, когда один и тот же образ доступен под несколько архитектур.
\end{itemize}

Все перечисленные объекты в составе образа называются \textit{блобами (blobs)}, идентифицируются дайджестом~--- криптографическим хешем (как правило, SHA-256) от их содержимого~--- и хранятся в registry независимо друг от друга. Связи между ними устанавливаются через ссылки, содержащие соответствующий дайджест. Такая контент-адресуемая модель обеспечивает целостность и неизменность образа: любое изменение содержимого приводит к изменению дайджеста, что делает образ ссылочно-целостным относительно своих компонентов.

\subsubsection{Слои образа}

Особый интерес для данной работы представляют слои образа, поскольку их совместное содержимое представляет файловую систему контейнера~--- то, что непосредственно анализируется сканером уязвимостей.

Слой представляет собой tar-архив, содержащий разностное содержимое файловой системы относительно предыдущего слоя в порядке наложения, определяемом манифестом. Иначе говоря, файловая система контейнера получается \textit{наложением} слоев друг на друга: каждый файл, имеющийся в более позднем слое, перекрывает одноименный файл из более раннего слоя, а специальный тип записи \textit{whiteout} позволяет логически удалить файл из более ранних слоев. Большинство реализаций container runtime используют для эффективного объединения слоев overlay-файловые системы, однако с точки зрения хранения и распространения каждый слой остается самостоятельным архивом.

Слои хранятся в registry в сжатом виде с использованием алгоритма gzip~\cite{gzip}. Уплотнение содержимого позволяет существенно сократить объем хранилища и сетевой трафик, поскольку файлы в составе образа в значительной части хорошо сжимаемы (текстовые файлы, скрипты, конфигурации). В контексте идентификации блобов с использованием хешей различают два типа дайджестов: \textit{compressed digest} соответствует сжатому виду слоя и используется на стороне registry для адресации блобов, в то время как \textit{diff ID} соответствует исходному, несжатому tar-архиву и используется в конфигурации образа для определения порядка наложения слоев. Подобное разделение обусловлено тем, что compressed digest зависит от выбранного алгоритма сжатия и его параметров, в то время как diff ID характеризует исключительно содержимое файловой системы и тем самым допускает альтернативные варианты сжатия без потери логической идентичности образа.

\subsubsection{Слои и кеширование}

Одним из важных свойств слоистой структуры OCI-образа является возможность \textit{совместного использования слоев} между различными образами. Если два или более образа имеют одинаковые базовые слои, то при их хранении в registry эти слои фактически содержатся в одном экземпляре~--- благодаря тому, что они адресуются дайджестом своего содержимого. Это же свойство активно используется при загрузке образов: container runtime может кешировать уже загруженные слои локально и пропускать их при загрузке других образов, использующих те же слои.

Сборщики образов, в свою очередь, как правило, организуют слои образа таким образом, чтобы максимизировать вероятность их повторного использования. На практике это приводит к характерной структуре, в которой образ состоит из нескольких слоев следующих типов:

\begin{itemize}
\item базовые слои, содержащие Linux-окружение, как правило, представленные публичными образами популярных дистрибутивов (Ubuntu, Alpine, Debian, Fedora) и редко изменяющиеся;
\item слои с зависимостями приложения~--- пакетами дистрибутива и языковыми пакетами, изменяющиеся при обновлении версий зависимостей;
\item один или несколько финальных слоев с самим приложением~--- кодом и/или скомпилированными бинарными файлами, которые обновляются на каждый коммит в репозитории.
\end{itemize}

Такая структура отражает естественную иерархию частоты изменения компонентов и хорошо сочетается с типовыми CI/CD-процессами, в которых при каждой сборке обновляется как правило только самый верхний слой с приложением, а остальные слои остаются неизменными.

\subsection{Container registry и хранение образов}

\textit{Container registry}~--- это сервис, отвечающий за хранение и распространение OCI-образов. На стороне registry образы хранятся в виде набора блобов, ссылающихся друг на друга в соответствии с описанной выше структурой. Все взаимодействие клиентов с registry осуществляется через HTTP-API, описанный в спецификации распространения OCI~\cite{ocidistspec}. С практической точки зрения наиболее интересны следующие ручки:

\begin{itemize}
\item \texttt{GET /v2/<name>/manifests/<reference>}~--- получение манифеста по имени образа и ссылке (тегу или дайджесту);
\item \texttt{GET /v2/<name>/blobs/<digest>}~--- получение блоба по его дайджесту;
\item \texttt{PUT /v2/<name>/manifests/<reference>}~--- загрузка манифеста;
\item \texttt{POST /v2/<name>/blobs/uploads/}, \texttt{PATCH}, \texttt{PUT}~--- последовательность операций загрузки блоба.
\end{itemize}

Важной особенностью протокола является то, что блобы могут запрашиваться как целиком, так и \textit{по диапазонам байтов} с помощью стандартных HTTP Range-запросов~\cite{httprange}. Это свойство, заложенное в самом протоколе HTTP, делает возможным частичную загрузку слоев~--- если содержимое слоя позволяет осмысленно интерпретировать его без полной загрузки, клиент может ограничиться запросом конкретных байтовых диапазонов. Эта возможность лежит в основе ряда форматов ленивой загрузки слоев и, в частности, формата eStargz, рассматриваемого в данной работе.

\subsubsection{Аутентификация и токены}

Для аутентификации в registry, как правило, применяется bearer-токен. При обращении к закрытым ресурсам незашифрованный клиент получает HTTP-ответ \texttt{401 Unauthorized} вместе с заголовком \texttt{Www-Authenticate}, содержащим адрес authorization-сервера, имя сервиса и запрашиваемые скоупы. Клиент обращается к authorization-серверу, обычно с использованием Basic-аутентификации, и получает токен, который затем используется в последующих запросах в заголовке \texttt{Authorization: Bearer <token>}. Этот механизм является стандартным для большинства публичных и приватных registry, включая Docker Hub, Harbor, GitHub Container Registry, Alibaba Container Registry и др.

\subsubsection{Managed container registry}

Особую категорию registry-инсталляций составляют \textit{managed container registry}~--- сервисы, предоставляемые облачными провайдерами. Их использование, как правило, освобождает пользователя от необходимости поддерживать собственную инфраструктуру registry: облачный провайдер берет на себя хранение, доступность и часто дополнительные функции~--- например, репликацию между регионами, сканирование на уязвимости, проверку подписей и т.п. Среди публичных примеров можно выделить Amazon Elastic Container Registry, Google Artifact Registry, Azure Container Registry, Yandex Cloud Container Registry.

С точки зрения нагрузки managed container registry представляет особый интерес, поскольку через него проходит значительный объем образов: организация может одновременно использовать registry для десятков и сотен проектов, каждый из которых регулярно загружает в registry новые версии образов в рамках своих CI/CD-процессов. Это создает существенные нагрузки на хранилище и сеть, а также делает задачу эффективного сканирования образов особенно актуальной.

\subsubsection{Harbor}

Среди open-source решений, ориентированных на корпоративное использование, выделяется проект Harbor~\cite{harbor}~--- зрелый container registry, перешедший в статус \textit{graduated} в Cloud Native Computing Foundation. Помимо базовой функциональности хранения и распространения, Harbor предоставляет ролевую модель доступа, репликацию между registry, поддержку подписей образов, политики удержания, аудит, а также интеграцию со сканерами уязвимостей через стандартизированный \textit{Pluggable Scanner API}~\cite{harborscanadapter}. На практике Harbor широко используется как локальная альтернатива публичным managed-сервисам и часто становится точкой интеграции сканера со всем циклом поставки.

\subsection{Сканирование как часть жизненного цикла}

Сканирование уязвимостей может производиться в различных точках жизненного цикла образа. С точки зрения момента запуска сканирования принято выделять следующие сценарии:

\begin{itemize}
\item \textit{во время сборки (build-time scanning)}~--- сканер вызывается на собранный образ непосредственно в CI-конвейере, ещё до его загрузки в registry. Результаты сканирования при необходимости позволяют прервать пайплайн в случае обнаружения уязвимостей выше заданного уровня критичности;
\item \textit{при загрузке в registry (push-time scanning)}~--- сканирование производится автоматически на стороне registry в момент или сразу после загрузки нового образа. Результаты сохраняются в самом registry и могут использоваться для последующих политик доступа и развертывания;
\item \textit{в произвольный момент (on-demand scanning)}~--- сканирование запускается администратором или автоматически по расписанию. Этот режим важен при появлении новых записей в базах уязвимостей: ранее проверенные образы могут оказаться уязвимыми относительно новой информации, и их следует пересканировать.
\end{itemize}

Во всех сценариях, кроме первого, сканер не имеет прямого доступа к локальной файловой системе с распакованным образом: образ к моменту сканирования уже находится на стороне registry. Соответственно, сканирование требует получения образа из registry. Именно эта модель~--- сканирование уже загруженных в registry образов~--- является основным объектом изучения данной работы. Она характерна как для managed container registry, так и для корпоративных инсталляций Harbor с включенной интеграцией со сканером, и именно в ней наиболее ярко проявляются ограничения существующих решений.

\newpage
